\chapter{Sécurité et préparation du déploiement}

\section{Constat réseau}

Les informations disponibles indiquent qu'aucun pare-feu ne sépare actuellement
la plateforme envisagée du réseau des automates. Cette situation augmente le
risque : une erreur de configuration pourrait rendre un service accessible
depuis une zone qui n'en a pas besoin.

\begin{figure}[H]
  \centering
  \includegraphics[width=\textwidth]{deploiement-securite.pdf}
  \caption{Déploiement proposé et mesures à appliquer avant le raccordement}
  \label{fig:security}
\end{figure}

\section{Mesures intégrées au prototype}

\begin{itemize}
  \item les services du projet ne commandent aucun équipement industriel ;
  \item Node-RED, InfluxDB, Mosquitto et le notifier sont liés à l'interface
        locale dans l'environnement de développement ;
  \item l'éditeur Node-RED demande un identifiant et un mot de passe ;
  \item les secrets restent dans \texttt{.env}, avec des permissions de fichier
        limitées ;
  \item Mosquitto possède un exemple de configuration TLS avec comptes séparés
        et ACL ;
  \item le flux Node-RED ne contient pas le token InfluxDB ;
  \item les messages invalides sont isolés dans un journal de rejet.
\end{itemize}

\section{Mesures requises sur site}

Avant toute connexion à Litmus Edge, plusieurs actions sont nécessaires :

\begin{enumerate}
  \item placer la VM de supervision dans un VLAN défini avec le responsable du
        réseau industriel ;
  \item appliquer un pare-feu hôte avec une politique de refus par défaut ;
  \item autoriser uniquement les connexions nécessaires entre Litmus Edge, le
        courtier et les interfaces utilisateur ;
  \item activer TLS pour MQTT et utiliser un certificat délivré ou approuvé par
        l'entreprise ;
  \item créer un compte MQTT par client et limiter ses topics ;
  \item créer un token InfluxDB en lecture seule pour Grafana et un token
        d'écriture distinct pour Node-RED ;
  \item intégrer Grafana à l'identité d'entreprise si elle est disponible ;
  \item sauvegarder le volume InfluxDB et tester une restauration sur une machine
        distincte.
\end{enumerate}

\section{Politique MQTT}

Litmus Edge doit pouvoir publier dans son espace de topics. Node-RED doit lire
ces topics et publier uniquement les rejets et son état. Grafana n'a pas besoin
d'accéder au courtier. Cette séparation réduit les droits de chaque service.

\begin{lstlisting}[caption={Extrait d'ACL MQTT de production}]
user litmus_edge
topic write sgx/<site>/#

user node_red
topic read sgx/<site>/#
topic write sgx/system/dead-letter
topic write sgx/system/node-red/status
\end{lstlisting}

\section{Données et conservation}

Le volume annoncé, inférieur à 200 pièces par heure sur une ligne, reste
modeste. Le prototype conserve tous les événements et compresse les mesures
anciennes. L'entreprise n'a pas encore fourni de durée officielle de
conservation. Cette durée devra être définie avec la qualité et la sécurité des
systèmes d'information avant la production.

Les sauvegardes doivent couvrir le volume InfluxDB, les volumes Grafana utiles
et les fichiers de configuration versionnés. Une sauvegarde qui n'a jamais été
restaurée ne constitue pas une preuve de reprise.
